\documentclass[11pt]{article}

\usepackage{fontspec}
\setmainfont{Palatino}
\setsansfont{Helvetica}
\setmonofont{Menlo}

\usepackage{amsmath, amssymb}
\usepackage[margin=1in]{geometry}
\usepackage{setspace}
\onehalfspacing
\usepackage{float}
\usepackage{graphicx}
\usepackage{booktabs}
\usepackage{xcolor}
\definecolor{blue}{RGB}{0,114,178}
\usepackage{hyperref}
\hypersetup{colorlinks, citecolor=blue, linkcolor=blue, urlcolor=blue}

\begin{document}

\title{What Lowi Predicted: Policy Type and the Committee Processing Penalty\\in the Korean National Assembly}
\author{KNA Research Agents (AI-generated)\thanks{Autonomously drafted by AI research agents. Not peer-reviewed or fact-checked. Do not cite. Source: \texttt{kna-research-agents.com}}}
\date{\today}
\maketitle
\thispagestyle{empty}

\begin{center}
\textbf{Keywords:} legislative committees, bill survival, policy typology, Korean National Assembly, redistributive politics
\end{center}

\begin{abstract}
\noindent Why do some bills advance through legislative committees while others languish indefinitely? I argue that the redistributive character of legislation creates a structural processing penalty at the committee stage, a prediction implied by Lowi's (1964) typology of policy types but never tested where most legislation actually dies. Using 15,291 classified bills from the Korean National Assembly (20th--21st terms, 2016--2024), I find that livelihood legislation targeting workers, care recipients, and welfare populations faces a 9.3 percentage-point committee processing penalty after controlling for arrival timing, sponsor characteristics, bill seriousness, and committee fixed effects. The penalty operates within committees and within individual sponsors' bill portfolios, ruling out compositional explanations. Redistributive labor bills receive committee decisions at nearly half the rate of distributive labor bills, precisely as Lowi's framework predicts. Under divided government, the content penalty nearly triples. These findings demonstrate that legislative effectiveness depends on what legislators attempt, not only on who they are.
\end{abstract}

\bigskip
\setcounter{page}{0}
\clearpage

%% =========================================================
\section{Introduction}
%% =========================================================

When a legislative committee decides which bills to advance and which to let expire, does the substantive content of the legislation matter? Lowi (1964) predicted that different policy types; distributive, regulatory, and redistributive; generate fundamentally different political dynamics, with redistributive proposals provoking the most intense organized opposition. If this prediction holds at the committee stage, where the vast majority of proposed legislation succeeds or fails, then the content of a bill should systematically predict whether it receives any committee action at all. Yet there exists, to my knowledge, no study that tests Lowi's prediction at the committee processing level in any legislature. This lacuna may stem from a measurement problem: classifying bills by policy type at scale requires text-based methods that have only recently become feasible, and most legislative datasets do not record the fine-grained committee-stage outcomes needed to observe where bills actually die.

The Korean National Assembly (KNA) provides an unusually powerful setting for this test. Approximately 80 percent of failed bills in the 20th--21st Assemblies (2016--2024) died at a single chokepoint: placed on a standing committee agenda but never receiving a committee decision. This is death by inaction, not rejection. The sheer volume of passive committee attrition; over 25,000 bills across two Assembly terms; creates the variation needed to ask whether content predicts which bills survive. The KNA's bill-tracking infrastructure, which records proposal dates, committee referrals, subcommittee review status, committee decisions, and plenary votes, provides the stage-by-stage data that most legislatures lack.

I exploit this setting to test whether Lowi's redistributive-distributive distinction predicts committee processing outcomes. Using keyword classification of bill proposal texts, I identify bills whose primary beneficiaries are workers (labor), care recipients, welfare populations, and small business owners; categories that Korean political discourse collectively labels \textit{minsaeng} (민생, ``livelihood'') legislation. I then classify bills within each domain by Lowi type: redistributive (mandating costs on identifiable groups) or distributive (providing targeted benefits without imposing concentrated costs). This two-dimensional classification; policy domain crossed with Lowi type; enables both between-domain and within-domain tests of the content penalty.

The results are striking. Livelihood bills face a 9.3 percentage-point penalty in committee processing probability relative to non-livelihood legislation, controlling for arrival timing, cosponsor count, proposal text length, sponsor party affiliation, and committee fixed effects ($\beta = -0.423$, $p < 0.001$). The penalty is not an artifact of bill-pool composition: among 327 legislators who introduced at least five bills in each category, their own livelihood bills fared 11.9 percentage points worse than their own non-livelihood bills ($t = -10.06$, $p < 0.001$). The penalty is equally large among bills with long, detailed proposal texts as among short ones, ruling out the interpretation that position-taking inflates the livelihood bill pool with frivolous legislation.

The Lowi decomposition reveals the mechanism. Within the labor domain, redistributive bills; minimum wage mandates, occupational safety penalties; receive committee decisions at a rate of 22.8 percent, compared to 38.5 percent for distributive labor bills such as employment subsidies and job training grants. This 15.7 percentage-point gap maps directly onto Lowi's prediction: redistributive legislation mobilizes organized opposition on both sides, while distributive legislation creates beneficiaries without generating concentrated losers. The gap vanishes in the care and welfare domains, where even nominally redistributive bills face only diffuse opposition from taxpayers rather than organized counter-mobilization. This pattern constitutes a scope condition on Lowi's theory: the redistributive penalty materializes only when the losing side possesses the organizational capacity to oppose.

These findings carry a broader theoretical implication. The legislative effectiveness literature, anchored by Volden and Wiseman (2014), has measured effectiveness as an attribute of \textit{legislators}; some are more skilled, more networked, or more strategically positioned than others. I show that effectiveness is also conditioned by \textit{what the legislator attempts}. The same sponsor achieves different committee processing rates depending on whether the bill redistributes or distributes, even holding constant all observable legislator characteristics. This Lowi-Volden synthesis; that policy type conditions legislative effectiveness; reframes the dependent variable in the effectiveness literature from ``who succeeds?'' to ``who succeeds at what?''

A secondary finding extends the analysis to political regime. The content penalty nearly triples under divided government: the average marginal effect of livelihood status shifts from $-7.0$ percentage points under unified government to $-18.4$ percentage points under divided government ($\beta_{\text{interaction}} = -0.536$, SE $= 0.089$, $p < 0.001$). Non-livelihood bills barely change across regimes. The distributional cost of divided government falls almost entirely on redistributive legislation; a finding that connects the committee processing penalty to the broader literatures on policy drift (Hacker 2004) and the political economy of inequality (Bonica et al.\ 2013).

The paper proceeds as follows. Section~2 reviews the literatures on policy typology, legislative winnowing, and legislative effectiveness, and derives expectations. Section~3 describes the KNA data, the keyword classification procedure, and the identification strategy. Section~4 presents the main results, robustness tests, and the divided government interaction. Section~5 discusses theoretical implications and limitations. Section~6 concludes.


%% =========================================================
\section{Literature and Theory}
%% =========================================================

\subsection*{Lowi's Typology and Committee Processing}

Lowi (1964) argued that different types of public policy generate distinct patterns of political conflict. Distributive policies allocate targeted benefits to specific groups without imposing visible costs on others, producing logrolling coalitions and low-conflict politics. Redistributive policies transfer resources from identifiable winners to identifiable losers, generating organized opposition on both sides and high-conflict bargaining. Regulatory policies fall between these poles. Wilson (1980) extended Lowi's framework into a two-by-two matrix of concentrated versus diffuse benefits and costs, but the core prediction remains: redistributive proposals should face the most political friction because they mobilize organized opposition.

Despite six decades of theoretical development, this prediction has not been tested at the committee processing stage; the precise institutional location where most legislation succeeds or fails. Kim (2019) applies Wilson's typology to the KNA, finding that policy type predicts which bills receive public hearings. But public hearings are a procedural decision; committee processing outcomes; whether a bill receives any substantive decision at all; are the consequential gate. I extend Kim's insight from procedure to substance.

\subsection*{Winnowing and Committee Capacity}

The literature on legislative winnowing offers a competing explanation for differential processing rates. Krutz (2005) demonstrates that Congress engages in systematic ``winnowing'' of the bill queue, with committee capacity constraints; not strategic selection; determining which bills receive attention. Under bounded rationality, committees use simple heuristics (arrival order, salience cues) to triage an overwhelming bill queue. If winnowing is content-blind, as Krutz's framework implies, then Lowi's prediction should fail at the committee stage: redistributive and distributive bills should die at equal rates once capacity constraints are controlled.

The KNA's bill volume makes the winnowing framework particularly relevant. Bill introductions increased roughly sixfold from the 17th Assembly (2004--2008) to the 21st Assembly (2020--2024), far outpacing committee capacity. Arrival timing emerges as the strongest predictor of committee action in my data: bills introduced in the first year of an Assembly term receive decisions at 42.7 percent compared to 27.5 percent for bills introduced in the fourth year; a 15.2 percentage-point gradient that dwarfs most other predictors. Content-blind winnowing is real in the KNA. The question is whether content \textit{additionally} predicts processing outcomes after controlling for the capacity-driven triage that Krutz identified.

\subsection*{Legislative Effectiveness and Policy Content}

Volden and Wiseman (2014) construct the Legislative Effectiveness Score (LES) to measure how successfully individual legislators advance their bills through the legislative process. Their framework decomposes legislative activity into sequential stages; introduction, committee action, floor passage, enactment; and aggregates across a legislator's portfolio. The LES has been applied to the U.S.\ Congress and extended to state legislative chambers. A key feature is that it treats all bills in a legislator's portfolio symmetrically: a legislator who passes five distributive bills and one who passes five redistributive bills receive equal credit. This aggregation implicitly assumes that policy type does not condition a legislator's ability to advance legislation; an assumption that Lowi's theory directly contradicts. If redistributive bills face structural opposition that distributive bills do not, then the LES conflates two sources of variation: legislator skill and bill difficulty. I test this assumption by comparing the same legislator's processing rates across policy types.

Volden, Wiseman, and Wittmer (2016) provide the closest precedent for content-classified bill fate analysis. They show that bills addressing women's issues in the U.S.\ Congress have a 2 percent passage rate compared to 4 percent overall, with women's issues bills sponsored by women themselves faring even worse. The parallel with the KNA data is striking: labor bills in my sample enact at 2.7 percent versus 7.0 percent for safety bills, yielding a comparable ratio. However, Volden, Wiseman, and Wittmer do not decompose the penalty into committee-stage versus floor-stage components, test for regime-conditional variation, or address the position-taking confound through within-sponsor comparisons. I extend their approach on all three dimensions.

\subsection*{Position-Taking and Bill Introduction}

Mayhew (1974) identifies position-taking as one of three reelection-motivated activities alongside advertising and credit claiming. If legislators introduce livelihood bills; minimum wage increases, childcare expansions; primarily to signal policy positions rather than to achieve legislative outcomes, then the lower processing rate of livelihood bills may reflect bill-pool quality rather than committee content bias. Kang and Park (2025) document systematic ``waffling'' in the KNA, where legislators sponsor bills they subsequently vote against, confirming that bill introduction and legislative intent are not always equivalent. Park (2023) shows that voters reward political grandstanding while remaining uninformed about legislative effectiveness, creating incentives for symbolic bill introduction on voter-salient topics.

The position-taking concern is serious but testable. If position-taking inflates the livelihood bill pool with frivolous legislation, the processing penalty should concentrate among ``unserious'' bills; those with short proposal texts, minimum cosponsors, and sponsors who lack institutional access to the reviewing committee. If the penalty persists among serious, well-crafted bills and within individual sponsors' portfolios, position-taking cannot be the sole explanation.

\subsection*{The Korean Committee Context and Expectations}

The KNA's standing committees exercise substantial gatekeeping authority. Committee chairs possess discretion over agenda scheduling, and the subcommittee review process (\textit{beopryulan-simsasoui}) serves as an additional filter. Kim and Lee (2023) demonstrate that a bill sponsor's membership on the reviewing subcommittee significantly predicts bill passage, establishing that institutional access matters for legislative outcomes in the KNA. An, Park, and Lee (2025) identify sponsor-level characteristics (party affiliation, seniority, committee position) as predictors of bill passage in the 20th--21st Assemblies, and Yoon and Shin (2023) show that cosponsorship network centrality is associated with legislative success. These studies examine the sponsor side of legislative outcomes. Despite this growing literature on \textit{who} succeeds in the KNA, there is a lack of systematic evidence on whether \textit{what} legislators propose independently predicts committee processing. I fill this gap.

Lowi's framework, applied to the committee stage, generates three testable expectations:

\textit{H1: Livelihood bills face a committee processing penalty relative to non-livelihood bills, controlling for sponsor characteristics and committee capacity.}

\textit{H2: Within policy domains, redistributive bills face a larger processing penalty than distributive bills, and this gradient is larger in domains with organized opposition (labor) than in domains with diffuse opposition (care, welfare).}

\textit{H3: The content penalty is amplified under divided government, when inter-party conflict raises the political costs of advancing redistributive legislation.}


%% =========================================================
\section{Data and Method}
%% =========================================================

\subsection{Data}

I draw on bill-level records from the Korean National Assembly's legislative information system for the 20th (2016--2020) and 21st (2020--2024) Assembly terms. The dataset contains 50,003 member-sponsored bills (\textit{uiwon-beopryulan}), excluding government-sponsored and committee-sponsored legislation. Each record includes the bill's proposal date, referred committee, subcommittee review status, committee decision (if any), plenary vote result (if any), proposal reason text (\textit{je-an-i-yu}), cosponsor list, and final legislative status.

The unit of analysis is the individual bill. The binary dependent variable, \textsc{Decision}, equals one if the bill received any committee decision; approval, rejection, or incorporation into an alternative bill (\textit{daean-bangyeong-pyegi}); and zero if the bill expired on the committee agenda without action. Among the 50,003 bills, 36.8 percent received a committee decision; 63.2 percent died from committee inaction.

I classify bills by policy domain using keyword matching on proposal reason texts. Bills are assigned to one of six domains; Labor, Care, Welfare, Small Business, Industry, and Safety; based on the presence of domain-specific Korean-language keywords in the proposal text. The classifier achieves 33.8 percent coverage, identifying 16,924 bills with clear domain assignments. Of these, 15,291 appeared on a committee agenda and constitute the analytical sample. The remaining 33,079 unclassified bills have a committee decision rate of 36.2 percent, falling between the livelihood (31.0 percent) and non-livelihood classified (40.4 percent) rates. I treat the unclassified residual as a scope condition: the findings apply to legislation with identifiable beneficiary populations, not to all legislation.\footnote{The keyword-assisted topic model approach (Eshima, Imai, and Sasaki 2023) would expand coverage while providing probabilistic classifications. I identify this as a priority upgrade for the publication-ready version.}

Within each domain, I classify bills by Lowi type. Bills whose proposal texts contain more redistributive keywords (mandates, penalties, prohibitions, regulatory strengthening) than distributive keywords (subsidies, tax credits, exemptions, incentives) are coded as redistributive ($N = 4{,}355$); those with the reverse pattern are coded as distributive ($N = 5{,}304$). This classification is admittedly crude; the word \textit{jiwon} (support) appears in both genuine subsidies and mandated accommodations; but it enables the within-domain test of Lowi's prediction that no prior study has attempted.

The key independent variables are: \textsc{Minsaeng}, an indicator for livelihood bills (Labor, Care, Welfare, or Small Business); \textsc{Redistributive}, an indicator for redistributive Lowi type; \textsc{Divided}, an indicator for bills processed under divided government (the Yoon presidency, 2022--2024). Controls include months since Assembly start (arrival timing), log cosponsor count, log proposal text length (a seriousness proxy), and sponsor party (Democratic Party of Korea indicator).

\subsection{Identification Strategy}

The main specification is a logistic regression of committee decision on policy content indicators:

\begin{equation}
\Pr(\textsc{Decision}_i = 1) = \Lambda\!\left(\beta_1 \textsc{Minsaeng}_i + \beta_2 \textsc{Redist}_i + \mathbf{X}_i \boldsymbol{\gamma} + \delta_c + \epsilon_i\right)
\label{eq:main}
\end{equation}

\noindent where $\Lambda(\cdot)$ is the logistic function, $\mathbf{X}_i$ is a vector of controls, and $\delta_c$ are committee fixed effects. The interaction model extends Equation~\ref{eq:main}:

\begin{equation}
\Pr(\textsc{Decision}_i = 1) = \Lambda\!\left(\beta_1 \textsc{Minsaeng}_i + \beta_3 \textsc{Divided}_t + \beta_4 (\textsc{Minsaeng}_i \times \textsc{Divided}_t) + \mathbf{X}_i \boldsymbol{\gamma} + \delta_c + \epsilon_i\right)
\label{eq:interaction}
\end{equation}

Identification rests on within-committee variation: comparing livelihood and non-livelihood bills referred to the same committee, with controls for bill-level characteristics. The committee fixed effects absorb time-invariant differences in committee workload, institutional norms, and chair behavior. The within-sponsor robustness test provides a second identification source: comparing the same legislator's livelihood and non-livelihood bills, which controls for all time-invariant legislator characteristics including political skill, ambition, and network position.

Two threats to identification deserve explicit acknowledgment. First, the keyword classifier may introduce measurement error that correlates with the treatment. If the keywords that define livelihood also capture political contentiousness independent of policy domain, the classifier picks up contentiousness rather than content. The Lowi decomposition partly addresses this concern by separating redistributive from distributive bills within each domain. Second, committee chair party affiliation, which may vary across committees and affect agenda scheduling; is not observed in these data. If opposition-party chairs systematically deprioritize the ruling party's signature livelihood legislation, the content penalty would partly reflect partisan gatekeeping. This confound is testable with chair roster data available on the KNA website, and I identify it as the most important gap for future work.


%% =========================================================
\section{Results}
%% =========================================================

\subsection*{Descriptive Patterns}

Table~\ref{tab:domain} reports committee decision rates and enactment rates by policy domain. A clear gradient emerges: labor bills receive committee decisions at 28.5 percent, while small business bills receive decisions at 46.4 percent. The risk ratio comparing labor to safety is $28.5 / 40.5 = 0.70$: labor bills are 30 percent less likely to receive any committee decision than safety bills. The enactment risk ratio ($2.7 / 7.0 = 0.39$) is even starker. These raw differentials motivate the multivariate analysis.

\begin{table}[H]
\centering
\caption{Committee Decision Rates by Policy Domain, 20th--21st KNA}
\label{tab:domain}
\begin{tabular}{lccc}
\toprule
Policy Domain & $N$ (on agenda) & Decision Rate (\%) & Enactment Rate (\%) \\
\midrule
Labor          & 2,887  & 28.5 & 2.7 \\
Care           & 2,605  & 31.6 & 3.1 \\
Welfare        & 2,524  & 34.0 & 5.3 \\
Industry       & 1,304  & 40.2 & 5.9 \\
Safety         & 4,057  & 40.5 & 7.0 \\
Small Business & 2,093  & 46.4 & 6.7 \\
\bottomrule
\multicolumn{4}{l}{\footnotesize Note: Unit is the individual bill. Decision includes approval, rejection, or}\\
\multicolumn{4}{l}{\footnotesize incorporation into an alternative bill. Enactment requires plenary passage}\\
\multicolumn{4}{l}{\footnotesize and presidential promulgation.}\\
\end{tabular}
\end{table}

The within-committee gaps confirm that the processing differential is not an artifact of bill sorting into differentially productive committees. Across the eight largest committees processing livelihood bills, the livelihood penalty ranges from $-6.8$ percentage points (Political Affairs Committee) to $-16.8$ percentage points (Environment and Labor Committee). The largest gap appears where the most contentious redistributive legislation; minimum wage mandates, occupational safety penalties; concentrates. The smallest gap appears in a committee where overall processing rates are low for all bill types, compressing the differential.

\subsection*{Multivariate Evidence}

Table~\ref{tab:main} reports logistic regression estimates from four nested specifications. Model~1 includes only the livelihood indicator. Model~2 adds bill-level controls and committee fixed effects. Model~3 adds the Lowi type indicator, sponsor party, and divided government status. Model~4 adds the interaction between livelihood status and divided government.

\begin{table}[H]
\centering
\caption{Logistic Regression: Committee Decision on Livelihood Status}
\label{tab:main}
\begin{tabular}{lcccc}
\toprule
 & (1) & (2) & (3) & (4) \\
 & Baseline & Controls + FE & Full & Interaction \\
\midrule
Minsaeng (livelihood)  & $-0.470^{***}$ & $-0.418^{***}$ & $-0.423^{***}$ & $-0.318^{***}$ \\
                        & $(0.041)$      & $(0.047)$      & $(0.048)$      & $(0.058)$ \\[4pt]
Months since start      &                & $-0.019^{***}$ & $-0.019^{***}$ & $-0.019^{***}$ \\
                        &                & $(0.002)$      & $(0.002)$      & $(0.002)$ \\[4pt]
Log cosponsors          &                & $-0.010$       & $+0.014$       & $+0.016$ \\
                        &                & $(0.021)$      & $(0.021)$      & $(0.021)$ \\[4pt]
Log text length         &                & $+0.204^{***}$ & $+0.213^{***}$ & $+0.210^{***}$ \\
                        &                & $(0.018)$      & $(0.018)$      & $(0.018)$ \\[4pt]
Redistributive          &                &                & $-0.093^{**}$  & $-0.096^{**}$ \\
                        &                &                & $(0.040)$      & $(0.040)$ \\[4pt]
DPK sponsor             &                &                & $-0.125^{***}$ & $-0.121^{***}$ \\
                        &                &                & $(0.038)$      & $(0.038)$ \\[4pt]
Divided government      &                &                & $+0.001$       & $+0.237^{***}$ \\
                        &                &                & $(0.043)$      & $(0.063)$ \\[4pt]
Minsaeng $\times$ Divided &              &                &                & $-0.536^{***}$ \\
                        &                &                &                & $(0.089)$ \\
\midrule
Committee FE            & No             & Yes            & Yes            & Yes \\
$N$                     & 15,291         & 15,291         & 15,291         & 15,291 \\
Pseudo $R^2$            & 0.010          & 0.044          & 0.045          & 0.046 \\
\bottomrule
\multicolumn{5}{l}{\footnotesize $^{*}p<0.10$, $^{**}p<0.05$, $^{***}p<0.01$. Standard errors in parentheses.} \\
\multicolumn{5}{l}{\footnotesize Coefficients are log-odds from logistic regression.} \\
\end{tabular}
\end{table}

The livelihood coefficient is negative, large, and statistically significant ($p < 0.001$) in every specification. Adding bill-level controls and committee fixed effects in Model~2 attenuates the coefficient by 11 percent (from $-0.470$ to $-0.418$); substantial but far from elimination. The average marginal effect (AME) from Model~3 is $-9.3$ percentage points: livelihood bills are 9.3 percentage points less likely to receive any committee decision than non-livelihood bills, all else equal. In a universe where the baseline committee decision rate is 35--40 percent, this represents a roughly 25 percent reduction in processing probability.

Two confound variables that might plausibly explain the gap do not. Arrival timing works in the \textit{opposite} direction: livelihood bills arrive earlier on average (month 18.0 from Assembly start) than non-livelihood classified bills (month 18.6). Earlier arrival should help, not hurt, and controlling for timing actually \textit{increases} the livelihood penalty. Cosponsor count has no predictive power: the coefficient is substantively zero and statistically insignificant across all models ($p = 0.60$ to $0.88$). In the KNA context, where the 10-cosponsor minimum is easily met and cosponsorship signatures are cheaply accumulated, cosponsorship is noise rather than signal; a negative finding with implications for comparative legislative studies that rely on cosponsor count as a proxy for legislative seriousness.

One control does matter: proposal text length ($\beta = 0.204$, $p < 0.001$). Bills with longer, more detailed proposal texts are significantly more likely to receive committee decisions, consistent with text length serving as a seriousness proxy. Even after this control, the livelihood penalty persists at full magnitude.

The Pseudo-$R^2$ of 0.046 indicates the model explains less than 5 percent of variation in committee processing. Individual bill fates depend heavily on unobservable factors: committee chair preferences, informal inter-party negotiations, executive branch lobbying, and bill-specific political dynamics. Published studies predicting individual bill outcomes routinely report comparable model fit. The relevant benchmark is not total explained variance but whether the content coefficient is meaningfully different from zero; and at an AME of 9.3 percentage points, it is both statistically and substantively significant.

\subsection*{The Lowi Decomposition}

Table~\ref{tab:lowi} decomposes decision rates within each policy domain by Lowi type. Two patterns emerge that confirm and refine the theoretical prediction.

\begin{table}[H]
\centering
\caption{Committee Decision Rates by Policy Domain and Lowi Type}
\label{tab:lowi}
\begin{tabular}{lccccr}
\toprule
 & \multicolumn{2}{c}{Redistributive} & \multicolumn{2}{c}{Distributive} & \\
\cmidrule(lr){2-3}\cmidrule(lr){4-5}
Policy Domain & Rate (\%) & $N$ & Rate (\%) & $N$ & Gap (pp) \\
\midrule
Labor          & 22.8 & 899   & 38.5 & 889   & $-15.7$ \\
Small Business & 31.3 & 412   & 53.3 & 1,116 & $-22.0$ \\
Care           & 32.4 & 709   & 32.3 & 932   & $-0.1$ \\
Welfare        & 35.3 & 563   & 35.2 & 1,061 & $+0.1$ \\
\bottomrule
\multicolumn{6}{l}{\footnotesize Note: Redistributive and distributive classification based on keyword}\\
\multicolumn{6}{l}{\footnotesize balance in proposal reason texts. Bills with equal counts excluded.}\\
\end{tabular}
\end{table}

First, labor and small business display massive Lowi gradients. Redistributive labor bills; minimum wage mandates, occupational safety penalties, compulsory insurance requirements; receive committee decisions at 22.8 percent, barely half the rate of distributive labor bills such as employment subsidies and job training grants. The small business gradient is even larger ($-22.0$ percentage points). These patterns map directly onto Lowi's prediction: redistributive legislation generates organized opposition from both the beneficiary side (workers demanding more) and the cost-bearing side (employers resisting mandates), while distributive legislation creates beneficiaries without mobilizing opponents.

Second, care and welfare show \textit{no} Lowi gradient. Redistributive and distributive care bills process at statistically identical rates ($-0.1$ percentage points), as do welfare bills ($+0.1$ percentage points). This null finding is as theoretically informative as the labor result. The redistributive penalty materializes only where the ``losing'' side of redistribution is well organized. For labor bills, employer federations (\textit{Gyeongchonghyeop}, the Korea Employers Federation; \textit{Daehan Sanggonghoeiso}, the Korea Chamber of Commerce) provide concentrated opposition to minimum wage increases and occupational safety mandates. For care and welfare bills, the ``losers'' are diffuse taxpayers who cannot organize to block specific proposals. This pattern identifies a scope condition on Lowi's theory, consistent with Olson's (1965) logic of collective action: the redistributive processing penalty requires concentrated opposition, not merely the abstract presence of redistribution.

This decomposition also resolves the small business puzzle from Table~\ref{tab:domain}. Small business bills have the highest overall decision rate (46.4 percent) despite being classified as livelihood legislation. The Lowi decomposition reveals that small business legislation is predominantly distributive ($N = 1{,}116$ versus $N = 412$ redistributive), and its distributive bills process at 53.3 percent; the highest rate in the sample. What distinguishes small business legislation is not political organization per se but its overwhelmingly distributive character: targeted subsidies, tax breaks, and market protections that create beneficiaries without identifiable losers.

The regression confirms the Lowi prediction with controls. The redistributive coefficient in Model~3 is $\beta = -0.093$ ($p = 0.021$), corresponding to a $-2.0$ percentage-point AME after controlling for livelihood status and all other covariates. Redistributive bills face a processing penalty \textit{on top of} the livelihood penalty.

\subsection*{Ruling Out Position-Taking}

Three independent tests address the concern that position-taking inflates the livelihood bill pool with frivolous legislation (Mayhew 1974; Kang and Park 2025).

\textit{Within-sponsor comparison.} Among 327 legislators who introduced at least five livelihood and five non-livelihood bills, the mean within-legislator gap is $-11.9$ percentage points ($t = -10.06$, $p < 0.001$). Seventy-two percent of these legislators show a negative gap; their own livelihood bills fare worse than their own non-livelihood bills. This result cannot be explained by between-legislator composition because it compares the same legislator's bills across categories. If a prolific position-taker introduces both types frivolously, both should show equally low processing rates. They do not.

\textit{Seriousness proxy stratification.} The livelihood penalty is essentially identical across subgroups defined by bill characteristics at introduction: $-10.8$ percentage points among bills with long proposal texts versus $-10.7$ among short texts; $-10.6$ among bills with above-median cosponsors versus $-11.1$ among those with the minimum. If position-taking inflated the livelihood pool with low-quality bills, the penalty should concentrate in the ``unserious'' subgroup. It does not.

\textit{Prolific livelihood sponsors.} Legislators in the top quartile of livelihood bill share show a within-portfolio gap ($-12.5$ percentage points) nearly identical to the population-wide gap. Position-taking may inflate the overall bill pool, but it does not drive the \textit{differential} processing penalty between livelihood and non-livelihood legislation.

\subsection*{The Divided Government Amplification}

Model~4 in Table~\ref{tab:main} includes the livelihood-by-divided-government interaction, which is the largest and most statistically significant coefficient in the analysis ($\beta = -0.536$, SE $= 0.089$, $p < 0.001$). The conditional AMEs tell a clear story. Under unified government (the Moon Jae-in presidency, 2020--2022), livelihood bills face a $-7.0$ percentage-point penalty relative to non-livelihood bills. Under divided government (the Yoon Suk-yeol presidency, 2022--2024), this penalty expands to $-18.4$ percentage points; nearly triple. Non-livelihood bills barely change across regimes (decision rates of 41.6 percent and 40.5 percent, respectively).

The regime transition did not slow the Assembly uniformly. It selectively stalled redistributive legislation while leaving other legislation relatively unaffected. This finding extends the gridlock literature (Binder 2003; Tsebelis 2002), which predicts that divided government reduces legislative output overall, by showing that the reduction falls disproportionately on legislation that would alter the distributional status quo. An illustrative case: all 31 minimum wage bills introduced in the 21st Assembly; addressing coverage expansion, sector-specific differentiation, disability wage exemptions, enforcement mechanisms, and commission reform; received zero committee decisions. They were 31 substantively distinct proposals, not duplicates. All expired with the Assembly term.


%% =========================================================
\section{Discussion}
%% =========================================================

The central finding; that the substantive content of legislation predicts committee processing outcomes, controlling for sponsor characteristics and institutional context; speaks to a foundational prediction that political science has theorized for six decades but never tested at the committee stage. Lowi (1964) predicted that redistributive policies generate distinct political dynamics. I show that these dynamics manifest as measurable processing differentials at the precise institutional location where most legislation dies.

The Lowi-Volden synthesis frames the contribution most precisely. Volden and Wiseman (2014) conceive of legislative effectiveness as a property of legislators. The within-sponsor penalty ($-11.9$ percentage points, $t = -10.06$) demonstrates that the same legislator, with the same skills and the same network position, achieves systematically different outcomes depending on whether the bill redistributes or distributes. This does not invalidate the LES framework but adds a dimension it currently lacks: policy type as a moderator of legislative effectiveness. The implication is that the legislative effectiveness literature should decompose effectiveness by policy type rather than aggregating across a legislator's full portfolio.

The null Lowi gradient in care and welfare provides a theoretically consequential scope condition. The redistributive penalty materializes only in domains where organized opposition exists (labor, small business), not in domains where the losing side is diffuse (care, welfare). Whether the mechanism operates through Lowi's content channel; policy type determines political dynamics; or through Olson's (1965) organizational channel; concentrated interests block legislation while diffuse interests cannot; is not distinguishable with the present data. Disentangling these channels would require measures of lobbying intensity or interest group activity at the bill level, data that the KNA does not systematically record. I interpret the finding as consistent with both channels and note that they generate identical observable predictions in the absence of lobbying data.

The existing Korean literature on legislative outcomes has focused extensively on the sponsor side. Kim and Lee (2023), An, Park, and Lee (2025), Yoon and Shin (2023), and Ka (2025) all examine which \textit{legislators} succeed at passing bills. My analysis complements this work by demonstrating that bill content independently predicts processing outcomes, controlling for the sponsor-level characteristics these studies identify. The implication is that legislative productivity research in the KNA should incorporate both content-side and sponsor-side predictors.

Several limitations qualify the findings. The keyword classifier covers only 33.8 percent of bills, creating a sample selection concern. The unclassified residual (66.2 percent of bills) may contain ``hidden livelihood'' legislation that uses different vocabulary; if these hidden livelihood bills face no processing penalty, the true population-level penalty is smaller than 9.3 percentage points.\footnote{The classified bills' decision rate (31.0 percent livelihood, 40.4 percent non-livelihood) brackets the unclassified bills' rate (36.2 percent), suggesting the residual is a mixture rather than a distinct category.} Committee chair party affiliation remains the most important unobserved confound. The low Pseudo-$R^2$ (0.046) confirms that individual bill fates are inherently difficult to predict and that the content penalty, while robust, explains only a small share of total variation. Finally, the within-legislator position-taking concern (where a legislator introduces a ``serious'' regulatory bill and a ``symbolic'' minimum wage bill) cannot be fully addressed without bill-level measures of post-introduction sponsor effort, such as whether the sponsor formally requested committee scheduling.


%% =========================================================
\section{Conclusion}
%% =========================================================

This paper provides a test of Lowi's policy type prediction at the committee processing stage, using 15,291 bills from the Korean National Assembly. Livelihood legislation faces a 9.3 percentage-point processing penalty after controlling for arrival timing, sponsor characteristics, and committee fixed effects. The penalty operates within committees, within individual sponsors' portfolios, and across all observable bill quality levels, ruling out compositional and position-taking explanations. The Lowi decomposition reveals the mechanism: the penalty concentrates in domains with organized opposition (labor, small business) and vanishes in domains with diffuse opposition (care, welfare), consistent with Lowi's typological prediction and Olson's collective action logic.

The findings carry implications for three literatures. For legislative effectiveness scholarship, they demonstrate that the unit of analysis should be the bill, not only the legislator: the same sponsor achieves different processing rates depending on the policy type of the legislation. For comparative legislative studies, they provide bill-level evidence that Lowi's framework has measurable consequences at the committee stage where most legislation is decided. For the political economy of inequality, they reveal a mechanism by which the legislative process perpetuates distributional advantage: legislation that would alter the status quo faces structural friction that legislation preserving it does not.

The divided government amplification carries normative weight. The transition from unified to divided government in the 21st Assembly did not impose equal costs across the policy agenda. It selectively froze redistributive legislation while leaving distributive legislation largely unaffected. Whether this represents democratic accountability; divided government checking policy overreach; or institutional pathology; divided government blocking redistribution; depends on normative commitments that the data cannot adjudicate. What the data establish is that the cost is not uniform, and that the populations who bear it are those whose legislative advocates face the steepest institutional odds.

\bigskip
\noindent\textit{This working paper was generated by AI research agents as an experimental output. It has not been peer-reviewed or fact-checked. Do not cite or use in any academic, policy, or professional context.}


%% =========================================================
\section*{References}
%% =========================================================

\begin{description}

\item An, Sungje, Sunchun Park, and Dongkyu Lee. 2025. ``A Study on the Factors Influencing the Passage of Legislation in the 20th and 21st National Assembly: Focusing on Bill Sponsors.'' \textit{The Journal of Korean Policy Studies} 25 (1): 115--.

\item Binder, Sarah A. 2003. \textit{Stalemate: Causes and Consequences of Legislative Gridlock}. Washington, DC: Brookings Institution Press.

\item Bonica, Adam, Nolan McCarty, Keith T. Poole, and Howard Rosenthal. 2013. ``Why Hasn't Democracy Slowed Rising Inequality?'' \textit{Journal of Economic Perspectives} 27 (3): 103--124.

\item Cox, Gary W., and Mathew D. McCubbins. 2005. \textit{Setting the Agenda: Responsible Party Government in the U.S. House of Representatives}. New York: Cambridge University Press.

\item Eshima, Shusei, Kosuke Imai, and Tomoya Sasaki. 2023. ``Keyword-Assisted Topic Models.'' \textit{American Journal of Political Science} 67 (4).

\item Hacker, Jacob S. 2004. ``Privatizing Risk without Privatizing the Welfare State: The Hidden Politics of Social Policy Retrenchment in the United States.'' \textit{American Political Science Review} 98 (2): 243--260.

\item Ka, Sangoon. 2025. ``Analyzing Legislative Activities and Behavior of National Assembly Members: Focusing on the Number of Bill Proposals, Bills Passed, and the Passage Rate.'' \textit{Journal of Research Methodology} 10 (3).

\item Kang, Sin-Jae, and Jiyoung Park. 2025. ``Why Do Legislators Engage in Waffling? Evidence from the Korean National Assembly, 2004--2020.'' \textit{Journal of East Asian Studies} 25.

\item Kim, Eun-Kyung. 2019. ``Analysing the Public Hearing in the National Assembly.'' \textit{Korean Journal of Policy Analysis and Evaluation} 16 (4).

\item Kim, Yanghun, and Dongseong Lee. 2023. ``An Analysis of the Impact of Bill Initiators' Position in Subcommittees on the Passage of Bills.'' \textit{Korean Political Science Review} 57 (1).

\item Krutz, Glen S. 2005. ``Issues and Institutions: `Winnowing' in the U.S. Congress.'' \textit{American Journal of Political Science} 49 (2): 313--326.

\item Lowi, Theodore J. 1964. ``American Business, Public Policy, Case-Studies, and Political Theory.'' \textit{World Politics} 16 (4): 677--715.

\item Mayhew, David R. 1974. \textit{Congress: The Electoral Connection}. New Haven, CT: Yale University Press.

\item Olson, Mancur. 1965. \textit{The Logic of Collective Action: Public Goods and the Theory of Groups}. Cambridge, MA: Harvard University Press.

\item Park, Hyeon-Ho. 2023. ``Electoral Rewards for Political Grandstanding.'' \textit{State Politics \& Policy Quarterly} 23 (4).

\item Tsebelis, George. 2002. \textit{Veto Players: How Political Institutions Work}. Princeton, NJ: Princeton University Press.

\item Volden, Craig, and Alan E. Wiseman. 2014. \textit{Legislative Effectiveness in the United States Congress: The Lawmakers}. New York: Cambridge University Press.

\item Volden, Craig, Alan E. Wiseman, and Dana E. Wittmer. 2016. ``Women's Issues and Their Fates in the U.S. Congress.'' \textit{Political Science Research and Methods} 6 (4): 679--696.

\item Wilson, James Q. 1980. \textit{The Politics of Regulation}. New York: Basic Books.

\item Yoon, Joochul, and Heontae Shin. 2023. ``Effects of Sponsor's Network Influence on Bill Success in Co-Sponsorship Network.'' \textit{Korean Public Administration Review} 57 (3): 97--.

\end{description}

\end{document}
